\chapter{Marco Teórico y Fundamentos}
\label{chap:marco-teorico}

% ============================================================================
% PROPÓSITO ÚNICO: Establecer contexto Smart Energy y fundamentos teóricos
% GENERALES. Detalles técnicos van en capítulos 3-4-5.
% ============================================================================

\section{Introducción}
\label{sec:marco-intro}

Este capítulo establece las bases teóricas que sustentan la arquitectura IoT propuesta para infraestructuras de Advanced Metering Infrastructure (AMI) en redes eléctricas inteligentes. La transición hacia Smart Grids demanda sistemas de medición que trasciendan la mera captura de consumo energético, integrando capacidades de monitoreo bidireccional, procesamiento distribuido en tiempo real, y conformidad con estándares internacionales de interoperabilidad y ciberseguridad~\cite{SmartHomeEnergy2024,velasquezSmartGridsEmpowered2024}.

La revisión bibliográfica se estructura en cinco secciones temáticas que progresan desde fundamentos conceptuales hacia el estado del arte en arquitecturas IoT para AMI. La Sección~\ref{sec:marco-smart-energy} contextualiza el ecosistema de estandarización que define el marco regulatorio y técnico (NIST Framework, IEEE 2030.5, ISO/IEC 30141), esencial para garantizar interoperabilidad multi-vendor y trazabilidad de datos en despliegues utility-scale. La Sección~\ref{sec:marco-edge-cloud} analiza paradigmas de procesamiento distribuido (edge computing vs cloud computing), identificando trade-offs de latencia, escalabilidad y privacidad que determinan la ubicación óptima de capacidad computacional. La Sección~\ref{sec:marco-seguridad} revisa amenazas específicas de sistemas IoT energy-constrained y mecanismos de defensa en profundidad que mitigan superficies de ataque ampliadas.

La Sección~\ref{sec:marco-estado-arte} constituye el núcleo del análisis, examinando tecnologías de comunicación LPWAN, protocolos mesh para redes de campo, plataformas de edge computing, y la emergencia de agentes de inteligencia artificial en el borde con modelos de lenguaje cuantizados. Esta revisión culmina en la identificación explícita de cuatro brechas en la literatura que motivan directamente las decisiones arquitectónicas del Capítulo~3. Finalmente, la Sección~\ref{sec:marco-conclusiones} sintetiza hallazgos clave y establece la transición conceptual hacia la arquitectura propuesta.

\section{Contexto Smart Energy y la Infraestructura AMI}
\label{sec:marco-smart-energy}

Esta sección presenta el marco de estándares internacionales que guían el diseño de infraestructuras AMI interoperables, desde el modelo de referencia NIST para Smart Grid hasta los protocolos de comunicación mandatorios (IEEE 2030.5) y los frameworks de arquitectura IoT (ISO/IEC 30141).

\subsection{Arquitectura de Referencia Smart Grid NIST}
\label{sec:marco-nist-framework}

Las Smart Grids integran tecnologías de información y comunicación (TIC) para monitoreo, control y optimización en tiempo real del flujo eléctrico desde generación hasta consumo final~\cite{SmartHomeEnergy2024,velasquezSmartGridsEmpowered2024}. Este enfoque permite: integración masiva de energías renovables distribuidas (DER), gestión activa de la demanda (DSM), detección y auto-recuperación de fallas (self-healing), y participación activa de prosumidores.

El modelo de referencia NIST para Smart Grid (NIST Framework Release 4.0~\cite{NISTFramework2022}) define tres capas principales~\cite{alsuwaidiSecuringSmartGrid2024}:

\begin{enumerate}
\item \textbf{Power and Energy Layer}: Infraestructura física de generación, transmisión, distribución y almacenamiento.
\item \textbf{Communication Layer}: Redes de datos multi-protocolo (HAN, NAN, WAN) que transportan información de telemetría y comandos de control.
\item \textbf{Application Layer}: Sistemas de gestión de energía (EMS), gestión de distribución (DMS), gestión de demanda (DERMS), y analytics.
\end{enumerate}

La arquitectura AMI se compone típicamente de: medidores inteligentes instalados en puntos de consumo, concentradores/gateways que agregan datos de decenas o cientos de medidores, y head-end systems en centros de control que procesan típicamente 1-10 millones de registros diarios en redes de 100K-1M medidores~\cite{alsafranChallengesImplementingIoT2025}.

\subsection{IEEE 2030.5-2018 (Smart Energy Profile 2.0)}
\label{sec:marco-ieee-2030-5}

IEEE 2030.5, anteriormente conocido como ZigBee SEP 2.0, es el estándar ampliamente adoptado para interoperabilidad de dispositivos Smart Energy en América del Norte, mandatorio para programas de respuesta a la demanda (Demand Response) en California según Senate Bill 2030 y Hawaii Rule 14H~\cite{IEEERecommendedPractice,CaliforniaCPUC2023}. Define un modelo RESTful sobre HTTP/TLS para comunicación cliente-servidor entre dispositivos de campo (medidores, termostatos, inversores solares) y sistemas de gestión (DERMS, head-end systems)~\cite{tangResearchInteroperabilityIoT}.

\subsubsection{Arquitectura RESTful del Estándar}

IEEE 2030.5 estructura funcionalidades en Function Sets exponiendo recursos REST~\cite{IEEERecommendedPractice}: /dcap (descubrimiento), /tm (sincronización horaria), /edev (registro dispositivos), /mup (datos de medición), /mr (perfiles de carga), /msg (notificaciones), /dr (Demand Response), /fsa (QoS).

\subsubsection{Function Sets Implementados}

\textbf{Device Capability (DCAP)}: El cliente consulta /dcap para descubrir qué Function Sets implementa el servidor.

\textbf{End Device (ED)}: Registro de dispositivos con LFDI (Long Form Device Identifier) derivado de certificado X.509, proporcionando $2^{160}$ identificadores únicos, garantizando ausencia de colisiones incluso en despliegues globales~\cite{IEEERecommendedPractice}.

\textbf{Mirror Meter Reading (MMR)}: Publicación de lecturas de medición con granularidad configurable (típicamente 15 minutos). Datos codificados en formato OBIS (Object Identification System) según IEC 62056 DLMS/COSEM~\cite{IEC62056-2021}.

\subsection{ISO/IEC 30141:2024 - Marco de Interoperabilidad IoT}
\label{sec:marco-iso-iec-30141}

ISO/IEC 30141, publicado originalmente en 2018 y actualizado en 2024~\cite{ISOIEC30141v2024}, define un framework estandarizado para sistemas IoT mediante cuatro vistas complementarias (funcional, información, despliegue, operacional). Especifica componentes, interfaces y flujos de información, complementando ISO/IEC 29100 (Privacy Framework) e ISO/IEC 27001 (Security Management)~\cite{ISOIEC30141v2024,tangResearchInteroperabilityIoT}.

\subsubsection{Modelo de Capas Funcionales}

ISO/IEC 30141 define cuatro vistas complementarias, siendo la Vista Funcional la más relevante para esta tesis. Descompone el sistema IoT en 7 entidades funcionales (FE)~\cite{tangResearchInteroperabilityIoT}: Sensing/Actuation (adquisición de datos y control), Processing (transformación y agregación), Storage (persistencia de datos), Communication (transporte entre FEs), Security (autenticación y cifrado), Management (configuración y actualizaciones OTA), Application Support (APIs y workflows).

\subsubsection{Mapeo de Arquitectura Propuesta a ISO/IEC 30141}

La arquitectura propuesta implementa las 7 entidades funcionales requeridas por ISO/IEC 30141:2024, garantizando conformidad con el estándar: (1) Sensing FE mediante medidores DLMS/COSEM y nodos Thread ESP32-C6, (2) Communication FE con gateway multi-radio Thread/HaLow, (3) Processing FE edge en Raspberry Pi 4 con ThingsBoard Edge Rule Engine, (4) Processing FE cloud con microservicios distribuidos y Kafka streams, (5) Storage FE local con TimescaleDB 7 días retention, (6) Storage FE cloud con ThingsBoard PostgreSQL 5 años, (7) Security FE con Thread AES-128-CCM-8, WPA3-SAE y mTLS 1.3, (8) Management FE con OTA firmware updates, (9) Application Support FE con REST API y IEEE 2030.5 adapter.

\section{Fundamentos de Computación Distribuida}
\label{sec:marco-edge-cloud}

% ============================================================================
% SECCIÓN §2.3: COMPUTACIÓN DISTRIBUIDA (Edge vs Cloud)
% Integrada desde: 02MarcoTeorico_SECCION_2_3_NUEVA.tex
% Redactada: 28 de noviembre de 2025 - FASE 5
% ============================================================================

% ============================================================================
% NUEVA SECCIÓN §2.3: COMPUTACIÓN DISTRIBUIDA (Edge vs Cloud)
% Redactada: 28 de noviembre de 2025
% Autor: GitHub Copilot - FASE 5
% Datos extraídos de: DATOS_CUANTITATIVOS_TOP20_PAPERS.md
% ============================================================================

\subsection{Arquitecturas de Procesamiento Distribuido: Edge Computing vs Cloud Computing}
\label{sec:edge-vs-cloud}

La evolución de infraestructuras IoT a escala masiva (millones de dispositivos) ha generado debate arquitectónico fundamental sobre la ubicación óptima del procesamiento de datos: centralizado en cloud datacenters con recursos virtualmente ilimitados, o distribuido en edge gateways cercanos a los dispositivos con recursos locales limitados pero latencia ultra-baja~\cite{andriuloEdgeComputingCloud2024,minhEdgeComputingIoTEnabled2022}. Esta dicotomía no representa elección binaria sino espectro continuo de arquitecturas híbridas que balancean trade-offs de latencia, bandwidth, escalabilidad, privacidad y costo según requisitos específicos de aplicación~\cite{andriuloEdgeComputingCloud2024}.

Para sistemas Smart Energy con integración de Distributed Energy Resources (DER), la elección arquitectónica presenta implicaciones críticas: monitoreo de power quality requiere procesamiento de waveforms con latencia <10 ms (incompatible con round-trip cloud), mientras analytics históricos de consumo agregado (millones de registros) demandan escalabilidad horizontal imposible en gateways edge con recursos limitados~\cite{minhEdgeComputingIoTEnabled2022}. Esta subsección analiza las características, ventajas y limitaciones de ambos paradigmas, fundamentando la arquitectura híbrida propuesta en esta tesis que implementa edge computing local para procesamiento tiempo real y sincronización bidireccional con cloud para analytics avanzados y dashboards remotos.

\subsubsection{Cloud Computing: Centralización con Escalabilidad Ilimitada}

El paradigma cloud computing concentra capacidad computacional, almacenamiento y servicios de red en datacenters remotos hyperscale (AWS, Azure, GCP) con arquitectura multi-tenant virtualizada, ofreciendo elasticidad bajo demanda y economías de escala mediante consolidación de workloads~\cite{andriuloEdgeComputingCloud2024}. Para aplicaciones IoT, el modelo típico implementa ingestion pipeline donde dispositivos transmiten telemetría vía WAN (LTE, fiber) a cloud ingestion layer (AWS IoT Core, Azure IoT Hub), persistiendo datos en databases escalables (Amazon RDS, Azure Cosmos DB) y ejecutando analytics mediante servicios serverless (AWS Lambda, Azure Functions).

\textbf{Ventajas Arquitectónicas Cloud:}

\begin{enumerate}
\item \textbf{Escalabilidad horizontal ilimitada:} Capacidad de provisionar recursos computacionales y storage dinámicamente según demanda. Un despliegue AMI de 1 millón de medidores generando 96 lecturas/día (intervalo 15 min) produce $1M \times 96 \times 2.5\text{ KB} = 240$ GB/día de telemetría~\cite{alsafranChallengesImplementingIoT2025}. Cloud storage S3/Blob con lifecycle policies (hot → warm → cold → archive) maneja este volumen con TCO \$0.01-0.05/GB/mes vs \$5-15/GB/mes para storage on-premise con SLAs comparables.

\item \textbf{Servicios gestionados (managed services):} Eliminan overhead operacional de mantenimiento de infraestructura. Amazon RDS (PostgreSQL/TimescaleDB) provee backups automáticos, replicación multi-AZ, patching de seguridad, y scaling sin downtime. Azure IoT Hub ofrece device provisioning, identity management, telemetry routing y monitoring con modelo OpEx pay-per-use vs CapEx on-premise.

\item \textbf{Analytics avanzados y ML/AI:} Acceso a servicios de machine learning (AWS SageMaker, Azure ML) para entrenamiento de modelos predictivos con datasets masivos (years of historical data). Ejemplo: modelo LSTM para forecasting de consumo energético entrenado con 5 años × 1M medidores × 96 lecturas/día = 175 TB de datos históricos, impracticable en edge gateways con storage local limitado (256 GB SSD típico).

\item \textbf{Disaster recovery y alta disponibilidad:} Datacenters multi-región con replicación geo-distribuida garantizan SLAs 99.95-99.99\% uptime. Backup automático incremental con retention configurable (7 días, 30 días, 7 años) cumple requisitos regulatorios de auditoría energética.

\item \textbf{Integración ecosistema:} APIs RESTful/GraphQL para integración con sistemas enterprise existentes (ERP SAP, billing systems, SCADA), dashboards web/mobile responsivos (React/Angular), y exportación de datos hacia data warehouses (Snowflake, BigQuery) para business intelligence.
\end{enumerate}

\textbf{Limitaciones Críticas Cloud para Smart Energy:}

\begin{enumerate}
\item \textbf{Latencia WAN inaceptable para control tiempo real:} Review sistemático de Andriulo et al. (2024)~\cite{andriuloEdgeComputingCloud2024} reporta latencias típicas cloud computing en rango \textbf{50-200 ms} considerando: (1) uplink device → cloud gateway (20-80 ms vía LTE según cobertura), (2) WAN routing Internet público (10-50 ms según geografía), (3) cloud ingestion processing (10-30 ms), (4) downlink response (20-80 ms). Esta latencia total 50-200 ms es \textbf{incompatible} con requisitos IEEE 2030.5 DERControl (<100 ms) para comandos de actuación críticos (disconnect switches, load shedding durante emergencias grid)~\cite{IEEERecommendedPractice}. Durante eventos transitorios de power quality (voltage sags <50 ms, inrush currents <20 ms), procesamiento cloud genera respuesta tardía post-evento sin capacidad de mitigación preventiva.

\item \textbf{Consumo bandwidth excesivo:} Transmisión continua de waveforms sin filtrado local consume bandwidth WAN desproporcionadamente. Escenario: 1000 medidores transmitiendo waveforms 10 kHz × 2 bytes × 2 canales = 320 kbps/medidor → throughput agregado 320 Mbps. Con backhaul LTE Cat-4 (150 Mbps downlink, 50 Mbps uplink), el uplink se satura con solo 156 medidores simultáneos, requiriendo upgrades costosos a fiber o 5G mmWave.

\item \textbf{Dependencia conectividad WAN:} Durante outages WAN (tormentas, desastres naturales, ataques DDoS), sistemas cloud-only quedan completamente inoperativos sin capacidad de monitoreo local. Utilities requieren resiliencia operacional durante emergencias cuando visibilidad grid es más crítica.

\item \textbf{Privacidad y regulación datos:} Transmisión de consumption patterns granulares (15 min) a cloud público plantea concerns de privacidad GDPR/CCPA. Inferencia de hábitos de ocupación mediante análisis de load curves facilita profiling no autorizado. Regulaciones emergentes (EU Data Act 2023) requieren data minimization y processing local donde posible.

\item \textbf{Costos OPEX escalables con volumen:} Modelo pay-per-use cloud presenta costos lineales con número de dispositivos y volumen de datos. Análisis TCO 10 años para 10,000 medidores (§5.5): AWS IoT Core \$1/millón mensajes + RDS PostgreSQL db.r5.4xlarge \$3,500/mes + S3 storage 100 TB/año = OPEX total \$500K-700K, vs gateway edge con buffer local 7 días + sync selectiva cloud reduciendo tráfico 90\% = OPEX \$50K-100K.
\end{enumerate}

\subsubsection{Edge Computing: Procesamiento Localizado con Baja Latencia}

Edge computing distribuye capacidad computacional y storage hacia el perímetro de la red (edge gateways, IoT gateways, fog nodes), procesando datos localmente cerca de los dispositivos fuente antes de transmitir hacia cloud solo agregados, alarmas o eventos relevantes~\cite{andriuloEdgeComputingCloud2024,minhEdgeComputingIoTEnabled2022}. Esta arquitectura invierte el paradigma cloud-centric: en lugar de "collect everything, process centrally", implementa "process locally, transmit selectively".

\textbf{Ventajas Arquitectónicas Edge:}

\begin{enumerate}
\item \textbf{Latencia ultra-baja para aplicaciones tiempo real:} Andriulo et al. (2024)~\cite{andriuloEdgeComputingCloud2024} reportan latencias edge computing típicas en rango \textbf{1-10 ms} para procesamiento local sin round-trip WAN. Descomposición latencia edge: (1) uplink device → gateway (1-5 ms vía Thread/HaLow local), (2) processing local (2-5 ms según complejidad algoritmo), sin downlink WAN. Reducción 5-20× vs cloud (50-200 ms) habilita control tiempo real para DER: detección de voltage sag y activación de Dynamic Voltage Restorer (DVR) dentro de 1 ciclo AC (16.67 ms @ 60 Hz), previniendo daño equipment sensitive~\cite{minhEdgeComputingIoTEnabled2022}.

\item \textbf{Reducción bandwidth 70-90\%:} Edge gateway ejecuta analytics local (agregación temporal, compresión lossy de waveforms, detección de anomalías) transmitiendo hacia cloud solo: (1) agregados estadísticos horarios (min/max/avg/stdev) vs raw samples cada 15 min, (2) eventos power quality detectados (sag/swell/interruption) con metadata vs waveforms completos, (3) alarmas críticas (threshold violations) vs polling periódico. Andriulo et al.~\cite{andriuloEdgeComputingCloud2024} estiman reducción bandwidth 70-90\% en deployments IoT industrial con filtrado edge, validado por case studies utilities (PG\&E California, Enel Italia).

\item \textbf{Operación autónoma durante outages WAN:} Gateway edge con storage local (256 GB SSD) y rule engine local (ThingsBoard Edge) mantiene funcionalidad crítica durante desconexiones WAN: dashboards locales vía LAN, alertas locales (email/SMS via modem backup), buffering de telemetría con sincronización automática al restaurar conectividad. Resiliencia crítica para infraestructura energética donde disponibilidad 99.99\% es mandatoria.

\item \textbf{Privacidad por diseño (privacy-by-design):} Procesamiento local de consumption data evita transmisión de información granular sensible hacia cloud público. Edge gateway puede implementar anonymization (aggregation temporal, differential privacy) antes de upload cloud, cumpliendo GDPR Article 25 y California CCPA requirements. Utilities evitan liability de data breaches cloud.

\item \textbf{Optimización costos OPEX long-term:} Inversión CAPEX inicial en gateways edge (\$500-800/gateway) se amortiza en 18-24 meses vs costos OPEX cloud recurrentes. Análisis TCO 10 años (§5.5): gateway edge con sync selectiva reduce costos cloud storage/compute 80-90\%, resultando en savings \$400K-600K para deployment 10,000 medidores.
\end{enumerate}

\textbf{Limitaciones Edge Computing:}

\begin{enumerate}
\item \textbf{Recursos computacionales limitados:} Gateway típico (Raspberry Pi 5, 8 GB RAM, ARM Cortex-A76 quad-core) no puede ejecutar modelos ML/AI pesados (LSTM deep learning, transformer architectures) que requieren GPUs con 16-32 GB VRAM. Analytics avanzados (clustering millones de consumption profiles, forecasting multi-variable) permanecen dominio cloud.

\item \textbf{Escalabilidad horizontal compleja:} Agregar capacidad edge requiere deployment físico de nuevos gateways con instalación on-site, vs cloud scaling con clicks en portal web. Upgrades de firmware OTA para fleet de 1000+ gateways presentan desafío operacional (bandwidth, rollback en fallas).

\item \textbf{Single point of failure local:} Falla de gateway edge (hardware failure, power outage sin UPS) causa pérdida de visibilidad de 100-500 medidores asociados, vs cloud multi-AZ con failover automático. Requiere estrategias redundancia (dual-gateway, mesh backup) incrementando CAPEX.

\item \textbf{Mantenimiento distribuido:} 1000 gateways edge requieren 1000 puntos de mantenimiento físico (vs 1 cloud datacenter), incrementando OPEX operacional para troubleshooting, upgrades hardware, reemplazo de componentes fallidos.
\end{enumerate}

\subsubsection{Arquitecturas Híbridas: Fog Computing y Edge-Cloud Synergy}

Andriulo et al. (2024)~\cite{andriuloEdgeComputingCloud2024} destacan que arquitecturas híbridas combinando edge computing y cloud computing representan \textbf{solución óptima} para mayoría de aplicaciones IoT industriales, superando limitaciones de ambos paradigmas puros. El concepto \textbf{fog computing} (propuesto originalmente por Cisco 2012, estandarizado por OpenFog Consortium 2017) define jerarquía multi-tier de procesamiento distribuido:

\begin{itemize}
\item \textbf{Tier 1 - Mist Computing:} Procesamiento ultra-edge en dispositivos IoT mismos (microcontrollers con TinyML)
\item \textbf{Tier 2 - Edge Computing:} Procesamiento en gateways locales (Raspberry Pi, industrial PCs)
\item \textbf{Tier 3 - Fog Computing:} Procesamiento en edge servers regionales (mini-datacenters utilities)
\item \textbf{Tier 4 - Cloud Computing:} Procesamiento centralizado en datacenters hyperscale
\end{itemize}

\textbf{Principio de Diseño Híbrido:} Ejecutar procesamiento en tier más bajo (cercano a datos) que satisface requisitos funcionales y no-funcionales, escalando hacia tiers superiores solo cuando necesario. Aplicado a Smart Energy:

\begin{table}[H]
\centering
\caption{Distribución Óptima de Workloads Edge vs Cloud para Smart Energy}
\label{tab:edge-cloud-workload-distribution}
\small
\resizebox{\textwidth}{!}{%
\begin{tabular}{|l|l|l|l|}
\hline
\rowcolor{gray!20}
\textbf{Workload} & \textbf{Latencia Req.} & \textbf{Tier} & \textbf{Justificación} \\
\hline
DER control (disconnect) & <100 ms & Edge & Latencia cloud 50-200 ms insuficiente \\
\hline
Power quality monitoring & <10 ms & Edge & Detección sag/swell requiere <1 ciclo AC \\
\hline
Waveform compression & <50 ms & Edge & Reduce bandwidth 90\% antes upload \\
\hline
Anomaly detection local & <1 s & Edge & Alertas inmediatas sin WAN dependency \\
\hline
Dashboards locales & <2 s & Edge & Operación durante outages WAN \\
\hline
Consumption forecasting & 1-10 min & Cloud & Requiere ML models pesados + históricos \\
\hline
Analytics multi-utility & 10-60 min & Cloud & Agregación millones medidores cross-region \\
\hline
Billing generation & Daily batch & Cloud & Integration ERP/CRM enterprise systems \\
\hline
Long-term storage (5 años) & N/A & Cloud & S3/Glacier lifecycle policies cost-effective \\
\hline
\end{tabular}}
\end{table}

\subsubsection{Contexto Smart Grid: Electricidad Mundial y Edge Computing}

Minh et al. (2022)~\cite{minhEdgeComputingIoTEnabled2022} proporcionan contexto global que fundamenta urgencia de arquitecturas edge para smart grids. El survey comprehensive (119 citaciones) reporta forecast de consumo eléctrico mundial aumentando \textbf{~70\% para 2050} (de 25 TWh actual a 42 TWh), impulsado por: (1) electrificación transporte (EVs), (2) electrificación calefacción/cooling (heat pumps), (3) growth económico economías emergentes, (4) digitalización economy (datacenters).

\textbf{Desafío Service Response Time:} Minh et al.~\cite{minhEdgeComputingIoTEnabled2022} identifican que el challenge significativo para monitoring y controlling smart grids modernos es \textbf{service response time}. Grids tradicionales uni-direccionales toleraban response times minutos-horas (manual switching, operator intervention). Smart grids bidireccionales con penetración masiva de DER (solar PV, wind, battery storage) requieren response time \textbf{<100 ms para control estabilidad} durante transitorios:

\begin{itemize}
\item \textbf{Voltage regulation:} Inyección solar PV variable (clouds) causa voltage swings 5-10\% en <1 segundo, requiriendo activación automática OLTC (On-Load Tap Changer) o capacitor banks dentro 2-3 ciclos AC.
\item \textbf{Frequency stability:} Pérdida súbita de generation (wind turbine trip) causa frequency deviation >0.2 Hz en <500 ms, requiriendo fast frequency response (FFR) de battery storage con latencia <100 ms.
\item \textbf{Islanding detection:} Desconexión grid segment debe detectarse en <2 s para activar anti-islanding protection y evitar backfeed peligroso hacia líneas en mantenimiento.
\end{itemize}

\textbf{Arquitectura Edge Propuesta por Minh et al.:} Edge computing servers instalados en subestaciones procesan datos de smart meters, PMUs (Phasor Measurement Units), y renewable generation monitoring, ejecutando algoritmos control local (voltage control, frequency regulation, demand response) con latencia target \textbf{<10 ms} end-to-end. Cloud backend recibe solo agregados para analytics históricos y planning long-term.

\subsubsection{Análisis Cuantitativo: Reducción Latencia Edge vs Cloud}

La Tabla~\ref{tab:edge-cloud-latency-comparison} consolida métricas cuantitativas de latencia extraídas de literatura, estableciendo baseline numérico para evaluación arquitecturas.

\begin{table}[H]
\centering
\caption{Comparación Cuantitativa Latencia: Edge vs Cloud Computing}
\label{tab:edge-cloud-latency-comparison}
\small
\resizebox{\textwidth}{!}{%
\begin{tabular}{|l|l|l|l|}
\hline
\rowcolor{gray!20}
\textbf{Componente Latencia} & \textbf{Edge} & \textbf{Cloud} & \textbf{Referencia} \\
\hline
\textbf{Uplink device $\rightarrow$ processor} & 1-5 ms & 20-80 ms & \cite{andriuloEdgeComputingCloud2024} \\
 & (Thread/HaLow local) & (LTE WAN) & \\
\hline
\textbf{Processing time} & 2-5 ms & 10-30 ms & \cite{andriuloEdgeComputingCloud2024} \\
 & (local gateway) & (cloud ingestion) & \\
\hline
\textbf{WAN routing} & 0 ms & 10-50 ms & \cite{andriuloEdgeComputingCloud2024} \\
 & (no WAN hop) & (Internet público) & \\
\hline
\textbf{Downlink response} & 0 ms & 20-80 ms & \cite{andriuloEdgeComputingCloud2024} \\
 & (local) & (LTE WAN) & \\
\hline
\textbf{Latencia TOTAL} & \textcolor{green}{\textbf{1-10 ms}} & \textcolor{red}{\textbf{50-200 ms}} & \cite{andriuloEdgeComputingCloud2024} \\
\hline
\textbf{IEEE 2030.5 DER req.} & \textcolor{green}{\checkmark} & \textcolor{red}{\xmark} & \cite{IEEERecommendedPractice} \\
 & Cumple <100 ms & Excede >100 ms & \\
\hline
\textbf{Power quality (<10 ms)} & \textcolor{green}{\checkmark} & \textcolor{red}{\xmark} & \cite{minhEdgeComputingIoTEnabled2022} \\
 & Viable & No viable & \\
\hline
\end{tabular}}
\end{table}

\textbf{Factor de Mejora Latencia:} Edge computing reduce latencia end-to-end por factor \textbf{5-20×} vs cloud computing (1-10 ms vs 50-200 ms). Esta reducción no es incremental sino \textbf{cualitativa}, habilitando categorías completas de aplicaciones control tiempo real imposibles con arquitectura cloud-only.

\subsubsection{Bandwidth Optimization: Reducción 70-90\% con Edge Filtering}

Andriulo et al.~\cite{andriuloEdgeComputingCloud2024} reportan que arquitecturas híbridas edge-cloud optimizan consumo bandwidth mediante filtrado y agregación local, reduciendo tráfico WAN uplink en \textbf{70-90\%} vs transmisión raw data cloud-only. Análisis cuantitativo para caso smart metering:

\textbf{Escenario baseline cloud-only (sin edge filtering):}
\begin{itemize}
\item 1000 medidores transmitiendo readings cada 15 min: $1000 \times \frac{24 \times 60}{15} = 96,000$ mensajes/día
\item Payload por mensaje: 250 bytes (timestamp, device ID, 10 metrics, checksums)
\item Bandwidth diario: $96,000 \times 250\text{ bytes} = 24$ MB/día = 2.22 kbps promedio
\item \textbf{Agregado con waveforms:} Si subset 100 medidores envían waveforms 10 kHz × 2 bytes durante 1\% duty cycle (14.4 min/día), bandwidth adicional: $100 \times 320\text{ kbps} \times 0.01 = 320$ kbps promedio → \textbf{Total 322 kbps}
\end{itemize}

\textbf{Escenario optimizado edge-cloud (con edge filtering):}
\begin{itemize}
\item Edge gateway agrega 1000 medidores localmente, transmitiendo hacia cloud solo:
  \begin{itemize}
  \item Agregados horarios (min/max/avg): $1000 \times 24 = 24,000$ mensajes/día ($4\times$ reducción)
  \item Waveforms comprimidos (wavelet transform, factor 10:1): 320 kbps / 10 = 32 kbps
  \item Eventos power quality (solo detecciones): ~100 eventos/día × 1 KB = 100 KB/día
  \end{itemize}
\item Bandwidth diario cloud upload: 6 MB + 0.4 MB + 0.1 MB = 6.5 MB/día = 0.6 kbps + 32 kbps = \textbf{32.6 kbps}
\item \textbf{Reducción bandwidth: $(322 - 32.6) / 322 = 89.9\%$} → Validando claim literatura 70-90\%
\end{itemize}

\textbf{Implicaciones TCO:} Reducción 90\% bandwidth disminuye costos WAN (LTE data plan, fiber circuits) en \$5-15/mes/gateway, acumulando savings \$60-180/año/gateway. Para deployment 100 gateways, savings anuales \$6K-18K OPEX WAN.

\subsubsection{Justificación Arquitectura Híbrida de esta Tesis}

La arquitectura propuesta en Capítulos 3-4 implementa modelo híbrido edge-cloud fundamentado en análisis comparativo presentado:

\begin{enumerate}
\item \textbf{Edge gateway local:} Raspberry Pi 5 ejecutando ThingsBoard Edge + TimescaleDB local + Kafka broker procesa waveforms localmente con latencia <10 ms, cumpliendo requisitos power quality monitoring y DER control <100 ms~\cite{andriuloEdgeComputingCloud2024,minhEdgeComputingIoTEnabled2022}.

\item \textbf{Sincronización selectiva cloud:} ThingsBoard Cloud backend recibe solo agregados horarios, eventos power quality detectados, y alarmas críticas. Reducción bandwidth 85-90\% vs cloud-only, disminuyendo OPEX WAN y cloud storage costs~\cite{andriuloEdgeComputingCloud2024}.

\item \textbf{Resiliencia WAN outages:} Buffer local 7 días (256 GB SSD) + rule engine local mantiene operación autónoma durante desconexiones WAN prolongadas, sincronizando automáticamente al restaurar conectividad. Cumple requisitos utilities de disponibilidad 99.99\%.

\item \textbf{Analytics avanzados cloud:} Modelos ML forecasting (LSTM consumption prediction), clustering profiles (K-means millones usuarios), y BI dashboards (Grafana/Tableau) ejecutan en cloud con datasets históricos 5 años, leveraging escalabilidad horizontal cloud.

\item \textbf{TCO optimizado:} Inversión CAPEX gateway \$500-800 amortizada en 18-24 meses vs OPEX cloud recurrente. Análisis detallado TCO 10 años en §5.5 demuestra savings 32-35\% arquitectura híbrida vs cloud-only (AWS IoT/Azure IoT Hub).
\end{enumerate}

\textbf{Trade-off fundamental aceptado:} Arquitectura híbrida introduce complejidad operacional adicional (gestión fleet gateways edge, firmware updates OTA, monitoring distribuido) vs simplicidad cloud-only. Esta complejidad se justifica por: (1) Requisitos latencia no negociables (<10 ms power quality, <100 ms DER control), (2) Reducción OPEX 32-35\% long-term, (3) Resiliencia operacional durante emergencies grid cuando visibilidad es más crítica~\cite{andriuloEdgeComputingCloud2024,minhEdgeComputingIoTEnabled2022}.

% ============================================================================
% FIN SECCIÓN §2.3
% ============================================================================


\section{Fundamentos de Seguridad en Sistemas IoT}
\label{sec:marco-seguridad}

% TODO: NUEVO contenido (1000 palabras)
% Enfoque GENERAL (detalles específicos en Cap 3-4)
% Longitud estimada: 3-4 páginas

\subsection{Amenazas y Vectores de Ataque}
\label{sec:marco-seguridad-amenazas}

% TODO: NUEVO contenido (500 palabras)
% - Ataques a dispositivos: Physical tampering, firmware compromise
% - Ataques a red: Eavesdropping, man-in-the-middle, DoS
% - Ataques a servidor: SQL injection, credential stuffing
% - Casos reales: Mirai botnet, Stuxnet, Industroyer
% Figura: Árbol de ataques IoT

\subsection{Marco de Seguridad NIST Cybersecurity Framework}
\label{sec:marco-seguridad-nist}

% TODO: NUEVO contenido (700 palabras)
% - NIST CSF 5 funciones:
%   1. Identify: Asset inventory, risk assessment
%   2. Protect: Access control, encryption, secure boot
%   3. Detect: Anomaly detection, intrusion detection
%   4. Respond: Incident response plan, forensics
%   5. Recover: Backup, disaster recovery
% - Aplicación a AMI: Mapeo funciones NIST a arquitectura (referencia Cap 3)
% Figura: NIST CSF wheel

\subsection{Mecanismos Criptográficos}
\label{sec:marco-seguridad-cripto}

% TODO: NUEVO contenido (500 palabras)
% - Cifrado simétrico: AES-128/256 (throughput vs seguridad)
% - Cifrado asimétrico: RSA, ECC (key exchange, firma digital)
% - Hash: SHA-256 (integridad)
% - TLS/DTLS: Handshake, cipher suites
% - Nota: Detalles implementación en Cap 3 (DTLS Thread), Cap 4 (TLS MQTT)

\section{Estado del Arte: Arquitecturas IoT para AMI}
\label{sec:marco-estado-arte}

% TODO: NUEVO contenido (1500 palabras) - CRÍTICO
% Revisión literatura para justificar brecha
% Longitud estimada: 5-6 páginas

% ============================================================================
% SECCIÓN §2.5.1: TECNOLOGÍAS LPWAN PARA SMART METERING
% Contenido importado desde archivo modular
% ============================================================================
% ============================================================================
% NUEVA SECCIÓN §2.4.1: TECNOLOGÍAS LPWAN PARA SMART METERING
% Redactada: 28 de noviembre de 2025
% Autor: GitHub Copilot - FASE 5
% Datos extraídos de: DATOS_CUANTITATIVOS_TOP20_PAPERS.md
% ============================================================================

\subsection{Tecnologías LPWAN para Smart Metering}
\label{sec:lpwan-technologies}

Las tecnologías de red de área amplia de bajo consumo energético (LPWAN - Low-Power Wide-Area Network) representan una familia de protocolos inalámbricos diseñados para conectar dispositivos IoT que requieren largo alcance (1-15 km), bajo consumo energético (operación con batería >5 años), y throughput moderado (0.3-250 kbps)~\cite{chaudhariLPWANTechnologiesEmerging2020,ugwuanyiSurveyIoTDeveloping2021}. A diferencia de tecnologías de corto alcance como Thread o Zigbee (alcance <500 m), las LPWAN permiten despliegues a escala utility con menor densidad de infraestructura, reduciendo CAPEX de gateways y concentradores~\cite{zhengPerformancePowerConsumption2018}.

Esta subsección analiza tres tecnologías LPWAN representativas con diferentes trade-offs de diseño: \textbf{(1) LoRaWAN} (red propietaria, sub-1 GHz, máximo alcance pero bajo throughput), \textbf{(2) NB-IoT} (red celular licenciada, infraestructura existente pero OPEX recurrente), y \textbf{(3) Wi-Fi HaLow (IEEE 802.11ah)} (sub-1 GHz unlicensed, throughput 20-40× superior a LoRa/NB-IoT). El análisis comparativo fundamenta la decisión arquitectónica de esta tesis de utilizar Wi-Fi HaLow como tecnología de backhaul para transmisión de waveforms de calidad de energía, dado que LoRaWAN/NB-IoT presentan throughput insuficiente (<250 kbps) para soportar muestreo continuo de forma de onda a 10 kHz (requisito mínimo ~200 kbps)~\cite{chaudhariLPWANTechnologiesEmerging2020}.

\subsubsection{LoRaWAN: Tecnología de Largo Alcance con Throughput Limitado}

LoRaWAN (Long Range Wide Area Network) es un protocolo MAC construido sobre la modulación propietaria LoRa (Chirp Spread Spectrum) de Semtech, operando en bandas ISM sub-1 GHz (868 MHz Europa, 915 MHz Américas, 433 MHz Asia). La arquitectura típica incluye end-nodes (sensores), gateways LoRaWAN (concentradores multi-canal), y network server (gestión de red y aplicaciones)~\cite{chaudhariLPWANTechnologiesEmerging2020}.

\textbf{Características Técnicas Clave:}

\begin{itemize}
\item \textbf{Alcance:} 2-5 km entorno urbano, hasta 15 km rural con line-of-sight~\cite{ugwuanyiSurveyIoTDeveloping2021}. El mayor alcance entre LPWAN se logra mediante spreading factors (SF7-SF12) que sacrifican data rate por robustez.

\item \textbf{Throughput:} Máximo teórico 50 kbps con SF7 @ 125 kHz bandwidth, reduciendo a 0.3 kbps con SF12 @ 125 kHz~\cite{chaudhariLPWANTechnologiesEmerging2020}. El duty cycle regulatorio del 1\% (restricción FCC/ETSI para bandas ISM) limita tiempo de transmisión acumulado a 36 segundos/hora, resultando en throughput efectivo promedio ~5-10 kbps para aplicaciones con alta frecuencia de reportes.

\item \textbf{Payload:} Limitado a 51 bytes por mensaje según especificación LoRaWAN 1.0.x~\cite{ugwuanyiSurveyIoTDeveloping2021}, requiriendo fragmentación de payloads mayores con impacto en latencia y confiabilidad.

\item \textbf{Consumo Energético:} Establece línea base de eficiencia energética en comparación con NB-IoT. Estudios experimentales de Ugwuanyi et al.~\cite{ugwuanyiSurveyIoTDeveloping2021} demuestran que LoRaWAN consume significativamente menos energía en operaciones de network join y transmisión de mensajes uplink, validando su idoneidad para aplicaciones battery-powered con reportes esporádicos (<1 mensaje/hora).

\item \textbf{Latencia:} Típicamente 1-5 segundos debido a arquitectura ALOHA pura sin coordinación de canal. LoRaWAN Class A (usado por >95\% dispositivos) solo recibe downlinks en ventanas RX1/RX2 después de uplink, introduciendo latencia adicional para comandos bidireccionales.
\end{itemize}

\textbf{Análisis para Smart Metering:} LoRaWAN es adecuado para telemetría simple de smart meters con lecturas horarias/diarias (payload <50 bytes, frecuencia <1 msg/hora) donde throughput limitado no representa restricción operacional. Sin embargo, presenta tres limitaciones críticas para arquitecturas AMI avanzadas con integración de Distributed Energy Resources (DER):

\begin{enumerate}
\item \textbf{Throughput insuficiente para waveforms:} La transmisión de forma de onda de voltage/corriente a 10 kHz requiere throughput mínimo de $10\text{ kHz} \times 2\text{ bytes/sample} \times 8 = 160$ kbps sin considerar overhead de protocolo~\cite{chaudhariLPWANTechnologiesEmerging2020}. LoRaWAN (50 kbps máximo teórico, 5-10 kbps efectivo con duty cycle 1\%) es \textbf{insuficiente por factor 16-32×}.

\item \textbf{Latencia incompatible con control DER:} IEEE 2030.5 Function Set DERControl especifica latencia <100 ms para comandos de actuación (disconnect switches, load shedding)~\cite{IEEERecommendedPractice}. LoRaWAN Class A con latencia típica 1-10 s no satisface este requisito.

\item \textbf{Payload limitado:} 51 bytes son insuficientes para reportes complejos de power quality que incluyen voltage sag/swell events, harmonic distortion measurements (THD), y phase imbalance metrics, requiriendo múltiples transmisiones fragmentadas con incremento de latencia end-to-end.
\end{enumerate}

\subsubsection{NB-IoT: LPWAN Celular con Infraestructura Licenciada}

NB-IoT (Narrowband IoT) es un estándar 3GPP Release 13 (2016) optimizado para comunicaciones machine-type (MTC) sobre infraestructura LTE existente, operando en bandas licenciadas 700-900 MHz. Utiliza tecnología narrowband (180 kHz bandwidth) para maximizar cobertura y penetración indoor, con tres modos de despliegue: standalone (reemplaza GSM), guardband (bandas de guarda LTE), o in-band (dentro de carrier LTE)~\cite{ugwuanyiSurveyIoTDeveloping2021}.

\textbf{Características Técnicas Clave:}

\begin{itemize}
\item \textbf{Alcance:} 10-15 km leveraging torres celulares existentes, con Maximum Coupling Loss (MCL) de 164 dB (superior a LoRaWAN ~157 dB)~\cite{ugwuanyiSurveyIoTDeveloping2021}. La penetración indoor profunda permite cobertura de medidores en sótanos y estructuras metálicas donde LoRaWAN falla.

\item \textbf{Throughput:} Máximo 250 kbps downlink (single-tone, 15 kHz subcarrier spacing) y 250 kbps uplink con multi-tone~\cite{chaudhariLPWANTechnologiesEmerging2020}. Mediciones experimentales de Ugwuanyi et al.~\cite{ugwuanyiSurveyIoTDeveloping2021} reportan throughput efectivo de 264 bps con latencia de 837 ms en testbed real, destacando el gap entre capacidad teórica y desempeño práctico en condiciones de cobertura marginal.

\item \textbf{Consumo Energético Comparativo:} Análisis experimental riguroso de Ugwuanyi et al.~\cite{ugwuanyiSurveyIoTDeveloping2021} con mismo hardware UE (User Equipment) para LoRaWAN y NB-IoT demuestra que NB-IoT consume significativamente más energía:
\begin{itemize}
    \item \textbf{Network Join:} NB-IoT consume +2.0 mAh adicionales vs LoRaWAN
    \item \textbf{Uplink Message (44 bytes):} NB-IoT consume +1.7 mAh adicionales por mensaje
    \item \textbf{Consumo Total Extra:} +3.7 mAh por operación típica (join + uplink)
\end{itemize}

Este overhead energético se debe a: (1) mayor complejidad de sincronización con red LTE (cell search, random access preamble, RRC connection setup), (2) potencia de transmisión más alta requerida para alcanzar MCL 164 dB, y (3) overhead de protocolo LTE (subheaders PDCP/RLC/MAC). Para aplicaciones battery-powered con transmisiones frecuentes (>10 mensajes/día), el consumo adicional acumulado de NB-IoT reduce significativamente la vida útil de batería vs LoRaWAN, aunque esta penalización se compensa parcialmente por Power Save Mode (PSM) y extended Discontinuous Reception (eDRX) que permiten duty cycles <0.01\%.

\item \textbf{Latencia:} 837 ms medido por Ugwuanyi et al.~\cite{ugwuanyiSurveyIoTDeveloping2021}, reducible a <100 ms con optimizaciones de red (early data transmission, control plane optimization). Latencia inferior a LoRaWAN pero superior a Wi-Fi HaLow.

\item \textbf{OPEX Recurrente:} A diferencia de LoRaWAN (banda ISM libre), NB-IoT requiere suscripción a operador celular con costo típico \$1-3 USD/dispositivo/mes en contratos enterprise~\cite{chaudhariLPWANTechnologiesEmerging2020}. Para despliegues de 10,000 medidores, el OPEX anual es \$120K-360K USD, representando 30-40\% del TCO a 10 años.
\end{itemize}

\textbf{Ventajas para Smart Metering:} NB-IoT presenta tres ventajas estratégicas sobre LoRaWAN: (1) \textbf{Cobertura garantizada} leveraging infraestructura celular existente de operadores, eliminando CAPEX de despliegue de gateways propietarios, (2) \textbf{Penetración indoor superior} con MCL 164 dB permite alcanzar medidores en ubicaciones desafiantes (sótanos, parking subterráneos, estructuras metálicas) donde LoRaWAN requeriría repetidores adicionales, (3) \textbf{Payload variable sin fragmentación} permite transmisión de mensajes >51 bytes en single uplink, reduciendo latencia end-to-end y simplificando lógica de reensamblado en application layer.

\textbf{Limitaciones para AMI Avanzado:} Similar a LoRaWAN, NB-IoT presenta throughput insuficiente (250 kbps teórico, 264 bps efectivo en cobertura marginal) para transmisión continua de waveforms de power quality a 10 kHz. La ecuación de Shannon-Hartley $C = B \log_2(1 + \text{SNR})$ aplicada a narrowband 180 kHz con SNR típico -6 dB (Extended Coverage Level 2) resulta en capacidad teórica máxima de ~50-100 kbps, insuficiente para 160 kbps requeridos por waveforms con overhead mínimo de protocolo~\cite{chaudhariLPWANTechnologiesEmerging2020}. Adicionalmente, el OPEX recurrente de \$1-3/mes/dispositivo impacta significativamente el TCO a 10 años vs tecnologías ISM unlicensed.

\subsubsection{Wi-Fi HaLow (IEEE 802.11ah): LPWAN de Alto Throughput para Smart Grid}

Wi-Fi HaLow es la denominación comercial del estándar IEEE 802.11ah (ratificado 2016), diseñado específicamente para IoT con operación en bandas sub-1 GHz (902-928 MHz US, 863-868 MHz EU) combinando las ventajas de propagación sub-GHz con throughput superior a tecnologías LPWAN tradicionales~\cite{zhengPerformancePowerConsumption2018,IEEE802.11ah-2016}. A diferencia de LoRaWAN/NB-IoT optimizados para narrow-band low-rate, 802.11ah adopta OFDM con ancho de banda configurable (1/2/4/8/16 MHz) permitiendo trade-off dinámico alcance-throughput.

\textbf{Especificaciones Técnicas IEEE 802.11ah:}

\begin{itemize}
\item \textbf{Frecuencia:} Bandas ISM sub-1 GHz unlicensed. US 902-928 MHz (26 MHz disponible), Europa 863-868 MHz (5 MHz disponible). Propagación superior a Wi-Fi 2.4/5 GHz por free-space path loss $\sim$15-20 dB menor y atenuación de paredes de concreto 8-12 dB menor según mediciones empíricas~\cite{zhengPerformancePowerConsumption2018}.

\item \textbf{Alcance:} Zheng et al.~\cite{zhengPerformancePowerConsumption2018} reportan ~1 km outdoor en entornos smart grid, validado por despliegue real de 600 millones de smart meters candidatos en China (2018). Análisis de link budget con potencia Tx +20 dBm, sensibilidad Rx -101 dBm (MCS0 @ 1 MHz), y path loss urbano COST231-Hata a 900 MHz resulta en alcance teórico 1.2-1.5 km, consistente con mediciones experimentales. Alcance 10-15× inferior a LoRaWAN pero suficiente para cobertura neighborhood-scale con densidad típica de medidores urbanos (200-500 medidores/km²).

\item \textbf{Throughput:} Rango 150 kbps (MCS0, 1 MHz, más robusto) hasta 86.7 Mbps (MCS10, 16 MHz, 4 spatial streams)~\cite{IEEE802.11ah-2016}. Para aplicaciones smart metering, configuración típica MCS7 @ 2 MHz entrega throughput efectivo ~4 Mbps con margen robusto para interferencia. \textbf{Análisis crítico para tesis:} 4 Mbps representa throughput 20-40× superior a LoRaWAN (50 kbps) y NB-IoT (250 kbps), permitiendo transmisión simultánea de waveforms de múltiples medidores: $\frac{4000\text{ kbps}}{200\text{ kbps/waveform}} = 20$ medidores simultáneamente con overhead de protocolo.

\item \textbf{Latencia:} CSMA/CA con Enhanced Distributed Channel Access (EDCA) para QoS diferenciado. Latencia típica 10-50 ms según prioridad (Voice/Video <20 ms, Best-Effort <50 ms), cumpliendo requisito IEEE 2030.5 DERControl de <100 ms para comandos críticos de actuación~\cite{IEEERecommendedPractice}.

\item \textbf{Eficiencia Energética:} Target Wake Time (TWT) permite scheduling determinístico de transmisiones con duty cycles <0.1\%, reduciendo consumo promedio a ~5-10 µA idle (similar LoRaWAN). Hierarchical Association Identifier (AID) soporta hasta 8,191 STAs (stations) por AP (Access Point), con grouping para beacon multicasting eficiente~\cite{IEEE802.11ah-2016}.

\item \textbf{Seguridad:} WPA3-SAE (Simultaneous Authentication of Equals) con Perfect Forward Secrecy, resistencia a offline dictionary attacks, y protección contra downgrade attacks. Superior a seguridad LoRaWAN (AES-128 vulnerable a replay attacks sin proper frame counters) y comparable a NB-IoT LTE security~\cite{IEEE802.11ah-2016}.
\end{itemize}

\textbf{Validación Industrial Smart Grid:} Zheng et al.~\cite{zhengPerformancePowerConsumption2018} analizan desempeño y consumo energético de IEEE 802.11ah para smart grid, destacando que China desplegó aproximadamente 600 millones de smart meters para 2018, muchos de ellos candidatos para conectividad HaLow debido a balance favorable entre throughput y alcance. El estándar 802.11ah fue publicado en 2017, posicionándose como tecnología emergente para Advanced Metering Infrastructure con capacidades superiores a ZigBee (limitado a 250 kbps @ 2.4 GHz) y LoRaWAN (limitado a 50 kbps).

\textbf{Caso de Uso Tesis - Transmisión de Waveforms:} La aplicación crítica que diferencia Wi-Fi HaLow de LoRaWAN/NB-IoT es la transmisión continua de formas de onda de calidad de energía para análisis avanzado de power quality. Consideremos un escenario de monitoreo de voltage/corriente con sampling rate de 10 kHz (suficiente para análisis hasta 5 kHz de bandwidth según criterio Nyquist):

\begin{equation}
\text{Throughput mínimo} = 10\text{ kHz} \times 2\text{ bytes/sample} \times 2\text{ canales (V+I)} \times 8\text{ bits/byte} = 320\text{ kbps}
\end{equation}

Con overhead de protocolo (headers UDP/IPv6, timestamps, device ID) estimado en 25\%, el throughput requerido es:

\begin{equation}
\text{Throughput con overhead} = 320\text{ kbps} \times 1.25 = 400\text{ kbps}
\end{equation}

\textbf{Comparación capacidades tecnologías LPWAN:}
\begin{itemize}
\item \textbf{LoRaWAN (50 kbps):} 400/50 = \textbf{8× insuficiente} → No viable para waveforms
\item \textbf{NB-IoT (250 kbps):} 400/250 = \textbf{1.6× insuficiente} → Marginal sin overhead, insuficiente con overhead real
\item \textbf{Wi-Fi HaLow (4 Mbps @ MCS7):} 4000/400 = \textbf{10× margen} → Viable con capacidad para 10 medidores simultáneos
\end{itemize}

Este análisis cuantitativo fundamenta la decisión arquitectónica de esta tesis de seleccionar Wi-Fi HaLow como tecnología de backhaul para el gateway smart meter. LoRaWAN y NB-IoT permanecen adecuados para telemetría simple (lecturas horarias/diarias sin waveforms), pero presentan limitaciones fundamentales de throughput para aplicaciones AMI avanzadas con integración DER que requieren visibilidad de power quality en tiempo real~\cite{chaudhariLPWANTechnologiesEmerging2020,zhengPerformancePowerConsumption2018}.

\subsubsection{Análisis Comparativo Multi-Criterio de Tecnologías LPWAN}

La Tabla~\ref{tab:lpwan-comparative-analysis} consolida las métricas cuantitativas extraídas de la literatura revisada, proporcionando framework de decisión multi-criterio para selección de tecnología LPWAN según requisitos específicos de aplicación Smart Energy.

\begin{table}[H]
\centering
\caption{Comparación Cuantitativa de Tecnologías LPWAN para Smart Metering}
\label{tab:lpwan-comparative-analysis}
\small
\begin{tabular}{|p{3cm}|p{2.8cm}|p{2.8cm}|p{2.8cm}|p{2.8cm}|}
\hline
\rowcolor{blue!20}
\textbf{Criterio} & \textbf{LoRaWAN} & \textbf{NB-IoT} & \textbf{Wi-Fi HaLow} & \textbf{Referencia} \\
\hline
\textbf{Alcance (km)} & \textcolor{green}{\textbf{15 km}} & 10-15 km & 1 km & \cite{ugwuanyiSurveyIoTDeveloping2021} \\
\hline
\textbf{Throughput} & 50 kbps & 250 kbps & \textcolor{green}{\textbf{10 Mbps}} & \cite{chaudhariLPWANTechnologiesEmerging2020} \\
\hline
\textbf{Latencia} & 1-5 s & 837 ms & \textcolor{green}{\textbf{10-50 ms}} & \cite{ugwuanyiSurveyIoTDeveloping2021} \\
\hline
\textbf{Payload típico} & 51 bytes & Variable & Variable & \cite{ugwuanyiSurveyIoTDeveloping2021} \\
\hline
\textbf{Consumo energético} & \textcolor{green}{\textbf{Baseline}} & +3.7 mAh & ~5 µA idle & \cite{ugwuanyiSurveyIoTDeveloping2021} \\
\hline
\textbf{Infraestructura} & Gateway propietario & \textcolor{green}{\textbf{Red celular}} & Access Point & - \\
\hline
\textbf{OPEX (USD/año)} & \textcolor{green}{\textbf{\$0}} & \$12-36 & \$0 & Análisis mercado \\
\hline
\textbf{Waveforms (>1 Mbps)} & \textcolor{red}{\xmark} & \textcolor{red}{\xmark} & \textcolor{green}{\checkmark} & \cite{zhengPerformancePowerConsumption2018} \\
\hline
\textbf{IEEE 2030.5 DER (<100 ms)} & \textcolor{red}{\xmark} & \textcolor{orange}{$\triangle$} & \textcolor{green}{\checkmark} & \cite{IEEERecommendedPractice} \\
\hline
\textbf{Despliegue documentado} & Millones (utilities EU/US) & Millones (telecom carriers) & \textcolor{blue}{\textbf{600M meters China}} & \cite{zhengPerformancePowerConsumption2018} \\
\hline
\end{tabular}
\end{table}

\textbf{Leyenda símbolos:} \textcolor{green}{\checkmark} Cumple requisito, \textcolor{red}{\xmark} No cumple, \textcolor{orange}{$\triangle$} Cumple marginalmente, \textcolor{green}{\textbf{Bold}} Mejor valor de criterio.

\subsubsection{Recomendaciones de Selección Según Caso de Uso}

El análisis cuantitativo comparativo permite establecer tres perfiles de aplicación Smart Metering con tecnología LPWAN óptima:

\textbf{Perfil 1 - Telemetría Simple (lecturas horarias/diarias):}
\begin{itemize}
\item \textbf{Requisitos:} Payload <50 bytes, frecuencia <1 msg/hora, latencia <10 s aceptable
\item \textbf{Tecnología óptima:} \textbf{LoRaWAN} (alcance máximo 15 km, costo OPEX \$0, eficiencia energética superior)
\item \textbf{Trade-off aceptado:} Throughput/latencia limitados no impactan aplicación simple
\end{itemize}

\textbf{Perfil 2 - AMI Estándar con cobertura indoor profunda:}
\begin{itemize}
\item \textbf{Requisitos:} Cobertura sótanos/estructuras metálicas, payload variable, latencia <1 s
\item \textbf{Tecnología óptima:} \textbf{NB-IoT} (MCL 164 dB, infraestructura celular existente, throughput 250 kbps)
\item \textbf{Trade-off aceptado:} OPEX recurrente \$12-36/año justificado por cobertura garantizada
\end{itemize}

\textbf{Perfil 3 - AMI Avanzado con DER y Power Quality (Tesis):}
\begin{itemize}
\item \textbf{Requisitos:} Waveforms >1 Mbps, comandos DER <100 ms, payload variable >100 bytes
\item \textbf{Tecnología óptima:} \textbf{Wi-Fi HaLow} (throughput 4-10 Mbps, latencia 10-50 ms, QoS EDCA)
\item \textbf{Trade-off aceptado:} Alcance reducido 1 km compensado por mayor densidad de APs en despliegue urban/suburban
\end{itemize}

\textbf{Justificación arquitectónica de esta tesis:} La selección de Wi-Fi HaLow como tecnología de backhaul se fundamenta en el \textbf{requisito crítico no negociable} de transmisión de waveforms de calidad de energía con sampling rate 10 kHz (equivalente a 400 kbps con overhead). LoRaWAN (50 kbps) y NB-IoT (250 kbps) son \textbf{técnicamente insuficientes} por factor 8× y 1.6× respectivamente. Esta limitación física de throughput no puede mitigarse mediante optimizaciones de protocolo o compresión de datos, dado que las formas de onda deben preservar fidelidad espectral hasta 5 kHz (criterio Nyquist) para detectar eventos de power quality como voltage sags/swells, harmonic distortion (THD), y transient overvoltages~\cite{chaudhariLPWANTechnologiesEmerging2020,zhengPerformancePowerConsumption2018}.

El despliegue documentado de 600 millones de smart meters en China considerando tecnología 802.11ah valida la madurez industrial de HaLow y reduce riesgo tecnológico de adopción. La arquitectura propuesta no excluye LoRaWAN/NB-IoT para telemetría simple de nodos sin capacidad de waveforms, permitiendo modelo híbrido multi-tecnología según perfil de medidor (básico vs avanzado)~\cite{ugwuanyiSurveyIoTDeveloping2021}.

% ============================================================================
% FIN SECCIÓN §2.4.1
% ============================================================================


\subsection{Arquitecturas Mesh para AMI}
\label{sec:marco-estado-mesh}

% TODO: NUEVO contenido (700 palabras)
% - Wi-Fi mesh (802.11s):
%   * Ejemplo: Silver Spring Networks (deprecated)
%   * Problema: Coexistencia 2.4 GHz, consumo alto
%   * Referencias: [Cite: Gungor et al. 2011, "Smart grid communications"]
% - Zigbee mesh:
%   * Ejemplo: Itron OpenWay CENTRON
%   * Problema: Limitado IPv6, 250 kbps max
%   * Referencias: [Cite: Sauter 2011, "The 3 families of IEC 62591"]
% - Thread mesh (emergente):
%   * Ventaja: IPv6 nativo, auto-formación
%   * Estado: Pocas implementaciones comerciales AMI
%   * Gap: **AQUÍ ESTÁ TU CONTRIBUCIÓN**

\subsection{Edge Computing en Smart Grids}
\label{sec:marco-estado-edge}

% TODO: NUEVO contenido (600 palabras)
% - Referencias:
%   * [Cite: Zhang et al. 2018, "Edge computing for smart grid"]
%     → Arquitectura conceptual, sin implementación HaLow
%   * [Cite: Okay et al. 2020, "Fog computing for smart cities"]
%     → Usa LTE backhaul (caro, alta latencia)
%   * [Cite: Liu et al. 2021, "Edge intelligence for IoT"]
%     → AI local, pero sin validación experimental AMI
% - Gap identificado:
%   * Falta arquitectura Thread + HaLow + Edge validada
%   * Falta análisis TCO detallado Edge vs Cloud
%   * Falta validación latencia E2E < 250 ms

\subsection{Inteligencia Artificial en Edge para IoT}
\label{sec:marco-estado-ia-edge}

La integración de capacidades de inteligencia artificial en dispositivos edge representa una tendencia emergente en sistemas IoT, particularmente relevante para aplicaciones Smart Energy que requieren análisis en tiempo real y toma de decisiones autónoma. Esta subsección presenta los fundamentos teóricos de IA en edge, Model Context Protocol (MCP) y sistemas multi-agente aplicados a gestión de infraestructura IoT.

\subsubsection{Fundamentos de IA Edge: LLMs y Quantización}
\label{sec:marco-estado-ia-fundamentos}

Los Modelos de Lenguaje Grande (\textit{Large Language Models}, LLMs) han demostrado capacidades de razonamiento y generación de lenguaje natural que pueden aplicarse a la gestión inteligente de sistemas IoT. Sin embargo, su despliegue en dispositivos edge enfrenta limitaciones de recursos computacionales (CPU/RAM) y energía que requieren técnicas de optimización.

\textbf{Quantización de modelos:} Técnica que reduce la precisión numérica de pesos y activaciones del modelo desde FP32 (32 bits) a representaciones de menor precisión (INT8, INT4) mediante técnicas post-entrenamiento. La quantización Q4\_K\_M (4-bit mixed precision) reduce el footprint de memoria hasta 75\% con degradación de perplexidad <5\%, permitiendo ejecutar modelos de 7-13B parámetros en hardware con 4-8 GB RAM~\cite{dettmersGPTQAccuratePost2023}.

\textbf{Modelos optimizados para edge:} Familias de modelos compactos diseñados específicamente para inferencia en dispositivos resource-constrained incluyen: Llama 3.2 (1B-3B parámetros, optimizado para móviles), Phi-3 Mini (3.8B parámetros, Microsoft Research, 2.3 GB Q4), y Mistral 7B (7B parámetros con context window extendido 32K tokens). Estos modelos balancean capacidad de razonamiento con eficiencia computacional~\cite{touvronLLaMAOpenEfficient2023,jiangMistral7B2023}.

\textbf{Latencia y throughput:} La inferencia LLM en dispositivos edge típicamente alcanza 10-50 tokens/segundo en CPU (Raspberry Pi 4, ARM Cortex-A72 @ 1.5 GHz) y 100-500 tokens/segundo con aceleración GPU/NPU. La latencia TTFT (\textit{Time To First Token}) crítica para aplicaciones interactivas oscila entre 200-2000 ms dependiendo del tamaño del modelo y prompt~\cite{zhaoTurningAILargeScale2024}.

\subsubsection{Model Context Protocol (MCP)}
\label{sec:marco-estado-mcp}

Model Context Protocol es un estándar abierto desarrollado por Anthropic (2024) que define una arquitectura cliente-servidor para conectar aplicaciones con modelos de lenguaje, proporcionando acceso estructurado a herramientas, datos y contexto externo mediante protocolo JSON-RPC 2.0~\cite{anthropicModelContextProtocol2024}.

\textbf{Arquitectura MCP:} Compuesta por tres elementos: (1) \textit{MCP Server} expone recursos y herramientas al modelo (e.g., API ThingsBoard Edge para consultas de telemetría), (2) \textit{MCP Client} integra el modelo en la aplicación (e.g., dashboard analítico operador), (3) \textit{Protocol Layer} gestiona comunicación bidireccional vía stdio, SSE o WebSockets.

\textbf{Casos de uso en IoT:} MCP permite a LLMs ejecutar acciones estructuradas sobre sistemas IoT: consultar telemetría histórica, generar reportes de anomalías, configurar dispositivos, y orquestar respuestas a eventos. La estandarización del protocolo habilita portabilidad de workflows entre diferentes proveedores de modelos (OpenAI GPT, Anthropic Claude, modelos locales Ollama)~\cite{anthropicModelContextProtocol2024}.

\textbf{Ventajas vs APIs REST tradicionales:} MCP proporciona: (1) Negociación automática de capabilities entre cliente y servidor, (2) Schema validation mediante JSON Schema de inputs/outputs, (3) Streaming de respuestas para latencias reducidas, (4) Gestión de contexto multi-turno con state management.

\subsubsection{Sistemas Multi-Agente para Gestión de Redes IoT}
\label{sec:marco-estado-multiagente}

Los sistemas multi-agente (\textit{Multi-Agent Systems}, MAS) distribuyen tareas complejas de gestión entre agentes especializados que colaboran para lograr objetivos del sistema. Aplicados a infraestructura IoT, permiten gestión autónoma de red, optimización de recursos y auto-healing~\cite{woldemarianFrameworkSelfHealingSystem2018}.

\textbf{Arquitecturas multi-agente:} (1) \textit{Jerárquica}: agente coordinador delega subtareas a agentes especializados (e.g., agente de red, agente de energía, agente de seguridad), (2) \textit{Distribuida}: agentes peer-to-peer negocian mediante protocolos de consenso, (3) \textit{Híbrida}: combinación de coordinación central con autonomía local de agentes edge~\cite{macasReviewMultiAgentSystems2018}.

\textbf{Agentes especializados para AMI:} (1) \textit{Network Agent}: monitoreo de conectividad, optimización de routing mesh, failover automático WAN, (2) \textit{Energy Agent}: análisis de consumo, detección de anomalías (hurto/fraude), forecasting de demanda, (3) \textit{Security Agent}: detección de intrusiones, gestión de certificados, respuesta a incidentes, (4) \textit{Maintenance Agent}: mantenimiento predictivo, scheduling de actualizaciones OTA, diagnóstico de fallos~\cite{ayaUnleashingIntelligenceEdge2024}.

\textbf{Comunicación inter-agente:} Protocolos de comunicación incluyen: (1) FIPA-ACL (Agent Communication Language) para mensajes estructurados, (2) MQTT para event-driven messaging IoT, (3) gRPC para RPC de baja latencia entre agentes co-ubicados. La coordinación puede implementarse mediante: (a) Blackboard pattern (pizarra compartida), (b) Contract Net Protocol (negociación basada en subastas), (c) BDI (Belief-Desire-Intention) para agentes con razonamiento explícito~\cite{macasReviewMultiAgentSystems2018}.

\subsubsection{Limitaciones y Brechas de Investigación}
\label{sec:marco-estado-ia-limitaciones}

La literatura reciente (2023-2025) documenta implementaciones experimentales de IA en edge para Smart Grids, pero identifica brechas críticas en validación de sistemas integrados~\cite{ayaUnleashingIntelligenceEdge2024,huangDataProcessingEnhancement2025}:

\textbf{Brecha 1 - Integración hardware-constrained:} Estudios existentes evalúan IA edge en hardware de propósito general (servidores x86, GPUs NVIDIA) no representativo de gateways AMI reales (ARM SoCs, presupuestos <10W). Raspberry Pi 4 (ARM Cortex-A72, 4 GB RAM, 8W) representa punto de equilibrio costo-rendimiento sin validación exhaustiva para inferencia LLM en producción.

\textbf{Brecha 2 - Latencia real-time:} Aplicaciones críticas como Demand Response (IEEE 2030.5) requieren latencias decisión <100 ms, no alcanzadas por LLMs edge (típicamente 500-2000 ms TTFT). Arquitecturas híbridas con procesamiento heurístico local + análisis LLM offline no han sido validadas empíricamente para AMI.

\textbf{Brecha 3 - MCP en producción:} Model Context Protocol (2024) carece de implementaciones documentadas en sistemas IoT de producción. La viabilidad de MCP Servers en gateways resource-constrained (overhead memoria/CPU, latencia tool calling) no ha sido caracterizada experimentalmente.

\textbf{Brecha 4 - Sistemas multi-agente escalables:} Propuestas teóricas de MAS para Smart Grids no validan escalabilidad >100 agentes ni overhead de coordinación en redes mesh reales. La integración de agentes IA con protocolos de gestión existentes (LwM2M, TR-069) permanece sin resolver.

\textbf{Oportunidades de investigación:} (1) Arquitecturas híbridas heurísticas-LLM para decisiones real-time con explicabilidad offline, (2) Técnicas de compresión agresiva (3-bit quantization, pruning estructurado) para modelos <1 GB ejecutables en gateways AMI, (3) Protocolos MCP-lite optimizados para enlaces intermitentes y alta latencia WAN, (4) Frameworks multi-agente con coordinación descentralizada resiliente a particiones de red.

\subsection{Justificación de la Brecha y Contribución de esta Investigación}
\label{sec:marco-estado-brecha}

El análisis crítico del estado del arte en arquitecturas IoT para AMI revela cuatro brechas fundamentales que motivan directamente las decisiones de diseño de esta tesis:

\textbf{Brecha 1: Ausencia de HaLow en gateways edge.} Ningún trabajo de los revisados integra Wi-Fi HaLow (IEEE 802.11ah) como backhaul primario para Smart Energy con failover multi-WAN, a pesar de su ventaja comprobada de alcance (1-3 km) y throughput (40 Mbps) sobre alternativas LPWAN como LoRaWAN (latencia típica >1s, throughput <50 kbps). Esta ausencia limita capacidades de transmisión de waveforms de calidad de energía que requieren throughput mínimo de 160-200 kbps para muestreo continuo a 10 kHz~\cite{chaudhariLPWANTechnologiesEmerging2020}. La arquitectura propuesta en el Capítulo~3 implementa gateway multi-radio con HaLow como tecnología principal y LTE como respaldo, habilitando monitoreo avanzado de power quality con latencia <15 ms.

\textbf{Brecha 2: Conformidad limitada IEEE 2030.5 + ISO/IEC 30141.} Soluciones comerciales (Cisco IR829, Dell Edge Gateway) implementan protocolos propietarios sin conformidad completa con IEEE 2030.5 Function Sets mandatorios (DERControl, Mirror Meter Reading). Prototipos académicos cumplen parcialmente estándares pero no documentan las cuatro vistas de ISO/IEC 30141 (funcional, información, despliegue, operacional), dificultando replicabilidad y certificación. Esta tesis provee mapeo explícito de arquitectura propuesta a las 7 entidades funcionales de ISO/IEC 30141 y documenta implementación completa de IEEE 2030.5 RESTful Function Sets mediante adapter ThingsBoard-to-SEP2.

\textbf{Brecha 3: Validación experimental insuficiente.} La mayoría de trabajos revisados presentan diseño conceptual o simulaciones sin mediciones empíricas de latencia end-to-end bajo carga, Packet Delivery Ratio (PDR) en condiciones reales con interferencia, o caracterización cuantitativa de resiliencia durante desconexiones WAN prolongadas (>24 horas). Solo 3 de los 20 papers analizados reportan deployment en campo con datos operacionales. El Capítulo~6 de esta tesis presenta validación experimental exhaustiva con: latencia percentil 95 medida en 10,000 transacciones, PDR con 500 dispositivos concurrentes, y pruebas de failover WAN→LTE con sincronización edge-cloud post-reconexión.

\textbf{Brecha 4: Costo-efectividad en contexto Latinoamericano.} Soluciones comerciales presentan barreras CAPEX significativas (Cisco IR829: \$2,500-4,000 USD, Dell EG: \$1,200-2,000 USD) para economías emergentes, sin análisis comparativo de alternativas open-source con hardware commodity. Esta tesis implementa gateway edge sobre Raspberry Pi 4 (\$55-75 USD) + módulo HaLow MM6108 (\$40-60 USD proyectado volumen), logrando CAPEX <\$200 vs >\$1,200 soluciones comerciales, con análisis de TCO a 5 años que considera OPEX de conectividad, mantenimiento y consumo energético.

\section{Conclusiones del Capítulo}
\label{sec:marco-conclusiones}

Este capítulo presentó los fundamentos teóricos que sustentan la arquitectura propuesta, abarcando contexto Smart Energy, computación distribuida edge, seguridad IoT, estado del arte en arquitecturas AMI, e inteligencia artificial aplicada a dispositivos edge.

\textbf{Contexto Smart Energy:} La infraestructura AMI enfrenta retos de escalabilidad (millones de medidores), latencia (Demand Response <5 segundos), y economía (CAPEX <\$100/dispositivo). Arquitecturas cloud-only (AWS IoT Core, Azure IoT Hub) presentan costos OPEX prohibitivos (\$50-200/device/año para 100K+ dispositivos) y vulnerabilidad a desconexiones WAN.

\textbf{Edge computing:} La computación en el borde proporciona equilibrio entre cloud-only y on-premise mediante procesamiento local, agregación de datos y operación autónoma durante fallos de conectividad. ThingsBoard Edge, plataforma open-source, reduce costos OPEX a cero versus alternativas comerciales (AWS Greengrass \$2.20/device/mes, Azure IoT Edge \$1.50/device/mes).

\textbf{Seguridad IoT:} Marco NIST CSF (Identify-Protect-Detect-Respond-Recover) aplicado a AMI requiere: autenticación mutual TLS dispositivos-gateway, segregación de redes (Thread mesh aislada de WAN), monitoreo de anomalías (CEP local), y resiliencia multi-WAN (failover LTE <5 segundos).

\textbf{Estado del arte:} La revisión de literatura (230+ papers 2018-2025) no identificó implementaciones que integren simultáneamente: Thread 1.3 mesh, HaLow 802.11ah long-range, edge computing ThingsBoard, y LwM2M device management. Trabajos previos documentan subconjuntos (Thread+edge~\cite{shahryarSolutionSmartHome2020}, HaLow+AMI~\cite{sunHaLowBasedWirelessCommunication2019}), pero ninguno valida arquitectura completa con métricas de producción. El análisis del estado del arte en inteligencia artificial edge (§2.4.4) identifica cuatro brechas críticas: integración hardware-constrained, latencia real-time <100 ms, MCP en producción IoT, y escalabilidad de sistemas multi-agente >100 nodos. Estas brechas justifican la Línea 6 de trabajo futuro (Cap. 7), donde LLMs quantizados, Model Context Protocol y arquitecturas multi-agente se proponen como evolución natural de la arquitectura validada en esta tesis.

\textbf{Transición al Capítulo 3:} Con el marco teórico establecido, el Capítulo 3 documenta la implementación del Nivel 1 de la arquitectura: nodos IoT con stack Thread + LwM2M ejecutando en microcontrolador ESP32-C6, validando consumo energético, latencia mesh, y gestión remota mediante Leshan LwM2M server.


